% id: 202-19
% book: 베이직쎈 대수 · 12 수학적 귀납법
% section: 유형 126 수학적 귀납법; 등식의 증명
% figure: none
% answer: 
% answer_source: 
% variant_level: 0 (원본 전사)
% 이 파일은 build.mjs 가 items.json 에서 만든다. 고칠 때는 items.json 을 고친다.
\smash{\raisebox{1.5mm}{\dmkichul{베이직쎈 대수 202쪽 19번}}}
다음은 모든 자연수 $n$에 대하여\par\smallskip\dispeq{$1+2+2^2+\cdots+2^{n-1}=2^n-1$}\par\smallskip\noindent 이 성립함을 수학적 귀납법으로 증명한 것이다.\exprbox{\begin{minipage}{0.96\linewidth}\raggedright\lineskip=4pt {\small\bfseries 증명}\par\smallskip (i) $n=1$일 때, \quad (좌변)$=1$, (우변)$=2^1-1=1$\\ \hspace*{1.5em}따라서 주어진 등식이 성립한다.\\ (ii) $n=\fbox{㈎}$일 때 주어진 등식이 성립한다고 가정하면\\ \hspace*{3em}$1+2+2^2+\cdots+2^{k-1}=2^k-1$\\ \hspace*{1.5em}위의 식의 양변에 $\fbox{㈏}$을 더하면\\ \hspace*{3em}$1+2+2^2+\cdots+2^{k-1}+\fbox{㈏}=2^{k+1}-1$\\ \hspace*{1.5em}따라서 $n=k+1$일 때도 주어진 등식이 성립한다.\\ (i), (ii)에서 모든 자연수 $n$에 대하여 주어진 등식이 성립한다.\end{minipage}}위의 과정에서 ㈎, ㈏에 알맞은 것을 구하시오.
